\chapter{積分と数列・不等式}\label{章:積分と数列}

積分を使うと，項数の増えていく和を一つの面積として見ることができる。また，積分の中身を大小比較すれば，値を正確に求められなくても積分そのものの範囲が分かる。この二つを組み合わせると，数列の極限や，複雑な積の極限まで調べられる。

この章では，まず和の中に区分求積の形を見つける。その後，積分を評価する不等式と，指数を一つずつ変えた積分を結び付ける漸化式を考える。最後に現れるベータ積分は，それまで別々に扱った多項式の積分と三角関数の積分を，一つの式の中に並べるためのものである。

\section{和の中から区分求積の形を取り出す}\label{節:和の区分求積}

連続関数$f$について，区間$[a,b]$を$n$等分し，右端を使って長方形を作ると，その面積の和は
\[
 \frac{b-a}{n}\sum_{k=1}^{n}
 f\left(a+\frac{k(b-a)}{n}\right)
\]
である。$n$を大きくすれば，この和は$\int_a^b f(x)\dd x$へ近づく。とくに$[0,1]$では
\begin{equation}\label{式:単位区間の区分求積}
 \lim_{n\to\infty}\frac1n\sum_{k=1}^{n}f\left(\frac{k}{n}\right)
 =\int_0^1 f(x)\dd x
\end{equation}
となる。

和の極限でこの公式を使うときには，$\frac{k}{n}$が横の位置，$\frac1n$が長方形の幅であることを確認する。分母に$n$が現れるというだけでは足りない。和全体を，幅と高さの積の和へ書き直す必要がある。

\begin{example}{分母をそろえて長方形の幅を見つける}
次の極限を求めよ。
\[
 \lim_{n\to\infty}\sum_{k=1}^{n}\frac{n}{n^2+k^2}.
\]
\end{example}

各項の分母から$n^2$をくくると
\[
 \frac{n}{n^2+k^2}
 =\frac1n\cdot\cfrac{1}{1+\left(\cfrac{k}{n}\right)^2}
\]
となる。したがって，これは$f(x)=\frac1{1+x^2}$の右端を使う区分求積である。この関数は$[0,1]$で連続なので
\[
 \lim_{n\to\infty}\sum_{k=1}^{n}\frac{n}{n^2+k^2}
 =\int_0^1\frac{1}{1+x^2}\dd x.
\]
$x=\tan t$とおけば，$0\leqq t\leqq\frac{\pi}{4}$で$\dd x=(1+\tan^2t)\dd t$だから，右辺は$\frac{\pi}{4}$である。和の一項は$0$へ近づくが，項数も増えるため，和全体は$0$にはならない。

式\eqref{式:単位区間の区分求積}で，右端の代わりに左端や中点を使っても，同じ積分へ収束する。ただし，和に含まれる点の範囲が変わるなら，積分区間も確かめなければならない。例えば
\[
 \frac1n\sum_{k=1}^{2n}f\left(\frac{k}{n}\right)
\]
では，幅は$\frac1n$のままで，点は$0$から$2$まで並ぶ。これは$[0,2]$の区分求積である。項数が$2n$だからといって，幅を勝手に$\frac1{2n}$へ変えてはいけない。

\begin{example}{和の範囲から積分区間を決める}
次の極限を求めよ。
\[
 \lim_{n\to\infty}\sum_{k=n+1}^{2n}\frac1k.
\]
\end{example}

一項を$\frac1k=\frac1n\cfrac1{\cfrac{k}{n}}$と書く。$k=n+1,\ldots,2n$に対応する点は
$1+\frac1n,\ldots,2$であり，$[1,2]$を$n$等分した小区間の右端である。よって
\[
 \lim_{n\to\infty}\sum_{k=n+1}^{2n}\frac1k
 =\int_1^2\frac1x\dd x=\log2.
\]
ここでは$\frac1x$が$x=0$で定義されないことは問題にならない。実際に使った積分区間は$[1,2]$であり，その上で連続だからである。

有限個の端の項を付けたり取り除いたりしても，その差が$0$へ近づくなら極限は変わらない。例えば，連続関数$f$について
\[
 \frac1n\sum_{k=0}^{n}f\left(\frac{k}{n}\right)
 -\frac1n\sum_{k=1}^{n}f\left(\frac{k}{n}\right)
 =\frac{f(0)}n\longrightarrow0
\]
である。しかし，取り除いた項数が$n$とともに増える場合や，端の関数値が大きくなる場合には，この確認を省けない。区分求積の公式を使う範囲を先に決め，残りの差を別に評価するのが確実である。

%
\section{単調な関数を使って和をはさむ}\label{節:和と積分の比較}

区分求積では，極限を取って和と積分を結び付けた。単調性が分かる場合には，有限の$n$でも大小関係が得られる。$f$が$[1,\infty)$で正かつ単調減少なら，$k\leqq x\leqq k+1$において
\[
 f(k+1)\leqq f(x)\leqq f(k)
\]
である。長さ$1$の区間で積分すると
\[
 f(k+1)\leqq\int_k^{k+1}f(x)\dd x\leqq f(k).
\]
この不等式を足すと
\begin{equation}\label{式:減少関数の和比較}
 \int_1^{n+1}f(x)\dd x
 \leqq\sum_{k=1}^{n}f(k)
 \leqq f(1)+\int_1^{n}f(x)\dd x
\end{equation}
を得る。右側では，$k=2,\ldots,n$について$f(k)\leqq\int_{k-1}^{k}f(x)\dd x$を足し，最初の$f(1)$だけを別に残している。

図\ref{図:和と積分の比較}では，各長方形の高さを小区間の左端の値にしている。減少関数では長方形が曲線の上に出るので，長方形の面積の和が積分の上からの評価になる。高さを右端の値に変えれば，大小関係が逆になる。

\begin{figure}[htbp]
 \centering
 \begin{tikzpicture}[x=1.15cm,y=2.9cm,>=Stealth]
  \foreach \k in {1,2,3,4} {
   \pgfmathsetmacro{\h}{1/\k}
   \fill[blue!12] (\k,0) rectangle ({\k+1},\h);
   \draw[blue!65!black] (\k,0) rectangle ({\k+1},\h);
  }
  \draw[->] (0.55,0)--(5.45,0) node[right]{$x$};
  \draw[->] (0.7,0)--(0.7,1.25) node[above]{$y$};
  \draw[thick,domain=1:5,samples=90] plot (\x,{1/\x});
  \node[anchor=west] at (2,0.86) {$y=\frac1x$};
  \foreach \k in {1,2,3,4,5}
   \node[below] at (\k,0) {$\k$};
 \end{tikzpicture}
 \caption{減少関数の左端の値で作った長方形は，曲線の上に出る。}
 \label{図:和と積分の比較}
\end{figure}

\begin{example}{調和数の増え方を調べる}
$H_n=\sum_{k=1}^{n}\frac1k$とおく。$H_n\to\infty$であること，および
\[
 \lim_{n\to\infty}\frac{H_n}{\log n}=1
\]
を示せ。後の比は$n\geqq2$について考える。
\end{example}

式\eqref{式:減少関数の和比較}を$f(x)=\frac1x$に用いると
\[
 \log(n+1)\leqq H_n\leqq1+\log n
\]
となる。下側の$\log(n+1)$が限りなく大きくなるため，$H_n$も限りなく大きくなる。次に，正の$\log n$で割れば
\[
 \frac{\log(n+1)}{\log n}
 \leqq\frac{H_n}{\log n}
 \leqq1+\frac1{\log n}.
\]
左辺は
\[
 \frac{\log(n+1)}{\log n}
 =1+\cfrac{\log\left(1+\cfrac1n\right)}{\log n}
 \longrightarrow1
\]
であり，右辺も$1$へ近づく。はさみうちの原理から求める極限が得られる。

ここから$H_n-\log n\to0$と結論してはいけない。比が$1$へ近づくということは，差が元の大きさに比べて小さいことを表すだけである。実際，上の評価から分かるのは，差が$0$以上$1$以下の範囲にあることまでである。

比較の向きを図から毎回考える必要はない。増加関数なら$f(k)\leqq f(x)\leqq f(k+1)$，減少関数ならその逆，と小区間ごとの不等式を書けばよい。添字のずれを避けるためにも，最初から大きな和の不等式を暗記するより，小区間から組み立てる方が扱いやすい。

%
\section{積分の大小とコーシー・シュワルツの不等式}\label{節:積分の不等式}

$f(x)\leqq g(x)$が区間全体で成り立てば，$\int_a^b f(x)\dd x\leqq\int_a^b g(x)\dd x$である。この基本的な性質を，積や平方の形に使うと，関数同士の関係を表す不等式が得られる。

\begin{theorem}{積分に対するコーシー・シュワルツの不等式}
$a<b$とし，$f,g$は$[a,b]$で実数値を取る連続関数とする。このとき
\begin{equation}\label{式:積分のコーシーシュワルツ}
 \left(\int_a^b f(x)g(x)\dd x\right)^2
 \leqq
 \left(\int_a^b f(x)^2\dd x\right)
 \left(\int_a^b g(x)^2\dd x\right)
\end{equation}
が成り立つ。$g$が恒等的に$0$でなければ，等号が成り立つための必要十分条件は，ある定数$c$について区間全体で$f(x)=cg(x)$となることである。$g$が恒等的に$0$なら，どの$f$についても等号が成り立つ。
\end{theorem}

平方は負にならないので，任意の実数$t$について
\[
 0\leqq\int_a^b\{f(x)-tg(x)\}^2\dd x
\]
である。式を短くするため
\[
 A=\int_a^b f(x)^2\dd x,\qquad
 B=\int_a^b f(x)g(x)\dd x,\qquad
 C=\int_a^b g(x)^2\dd x
\]
とおくと，右辺は$A-2Bt+Ct^2$になる。

$C>0$のとき，$t=\frac BC$を代入すれば
\[
 0\leqq A-\frac{B^2}{C}
\]
だから，$B^2\leqq AC$が得られる。等号が成り立つ場合には
\[
 \int_a^b\left\{f(x)-\frac BCg(x)\right\}^2\dd x=0.
\]
被積分関数は連続かつ非負である。もし一点で正なら，連続性によってその近くの小区間でも正となり，積分が正になってしまう。したがって区間全体で$f(x)=\frac BCg(x)$である。逆に$f=cg$なら直接代入して等号を確認できる。

$C=0$の場合には，同じ連続性の議論によって$g$が区間全体で$0$となる。このとき$B=0$でもあるので，不等式の両辺はともに$0$である。$C$で割る計算の前にこの場合を分けたのは，$0$による除算を避けるためである。

この不等式を使うときには，評価したい積分の中身を$f(x)g(x)$という積に見る。片方を$1$としてよいことも多い。例えば$g(x)=1$とすれば
\begin{equation}\label{式:平均と二乗平均}
 \left(\int_a^b f(x)\dd x\right)^2
 \leqq(b-a)\int_a^b f(x)^2\dd x.
\end{equation}
区間の長さで割ると，関数の平均値の平方は，平方の平均値を超えないことが分かる。等号は$f$が定数関数のときに限る。

\begin{example}{正の関数とその逆数を組み合わせる}
$h$は$[a,b]$で連続で，すべての$x$について$h(x)>0$を満たすとする。
\[
 \left(\int_a^b h(x)\dd x\right)
 \left(\int_a^b\frac1{h(x)}\dd x\right)
 \geqq(b-a)^2
\]
を示し，等号条件を求めよ。
\end{example}

二つの積分をコーシー・シュワルツの不等式の右辺にそろえたい。そこで$f=\sqrt h$，$g=\frac1{\sqrt h}$とおけば，$f^2=h$，$g^2=\frac1h$となる。$h$は閉区間で連続かつ正なので，この二つの関数も連続である。さらに$fg=1$だから
\[
 (b-a)^2=\left(\int_a^b1\dd x\right)^2
 \leqq\left(\int_a^b h\dd x\right)
       \left(\int_a^b\frac1h\dd x\right).
\]
等号条件$f=cg$は$\sqrt h=\frac c{\sqrt h}$，すなわち$h=c$を意味する。したがって，等号は$h$が正の定数関数のとき，またそのときに限って成り立つ。

例えば$h(x)=1+x$，$0\leqq x\leqq1$なら，$h$は定数ではないから
\[
 \frac32\log2>1
\]
となる。ここでは対数の値を小数で計算せずに，$\log2>\frac23$という厳密な下からの評価が得られた。

%
\section{指数を変えた積分を結ぶ}\label{節:積分漸化式}

積分に自然数$n$が含まれているとき，毎回展開して計算するより，$n$を一つ小さくした積分と結び付ける方が簡単な場合がある。部分積分で$x^n$を微分すれば$n x^{n-1}$が現れるので，指数の異なる積分の間に関係を作れる。

\begin{example}{部分積分から漸化式を作る}
非負整数$n$に対して
\[
 A_n=\int_0^1 x^n\e^{-x}\dd x
\]
とおく。$A_n$の漸化式を求め，$A_2$を計算せよ。また$A_n\to0$を示せ。
\end{example}

$n\geqq1$のとき，$x^n$を微分し，$\e^{-x}$を積分する部分積分を行うと
\begin{align*}
 A_n
 &=\left[-x^n\e^{-x}\right]_0^1
   +n\int_0^1x^{n-1}\e^{-x}\dd x\\
 &=-\e^{-1}+nA_{n-1}.
\end{align*}
$x=0$の端点の項が$0$になるのは$n\geqq1$だからである。最初の値は別に
\[
 A_0=\int_0^1\e^{-x}\dd x=1-\e^{-1}
\]
と求める。よって
\[
 A_1=1-2\e^{-1},\qquad A_2=2-5\e^{-1}.
\]
一方，極限を調べるには，漸化式より積分の評価が短い。$[0,1]$では$0<\e^{-x}\leqq1$なので
\[
 0<A_n\leqq\int_0^1x^n\dd x=\frac1{n+1}\longrightarrow0.
\]
漸化式は値を順に求めるのに役立ち，積分表示は符号や大きさを調べるのに役立つ。同じ数列でも，問いに応じて使いやすい表し方が違う。

漸化式を一般項まで解くこともできる。$A_n=nA_{n-1}-\e^{-1}$を$n!$で割ると
\[
 \frac{A_n}{n!}=\frac{A_{n-1}}{(n-1)!}-\frac{\e^{-1}}{n!}.
\]
$1$から$n$まで足せば中間の項が消え，$0!=1$として
\[
 A_n=n!\left(1-\e^{-1}\sum_{k=0}^{n}\frac1{k!}\right)
\]
を得る。ただし，この形だけを見て$n\to\infty$を取ろうとすると，括弧の中がどれほど速く$0$に近づくかも調べる必要がある。先ほどの積分による評価は，その難しさを避けて直接$A_n\to0$を示している。

%
\section{三角関数のべきとウォリスの公式}\label{節:ウォリス公式}

次に
\[
 I_n=\int_0^{\frac{\pi}{2}}\sin^n x\dd x
 \qquad(n=0,1,2,\ldots)
\]
を考える。最初の二つは
\[
 I_0=\frac{\pi}{2},\qquad I_1=1
\]
である。$\sin^n x$を$\sin^{n-1}x\sin x$と分け，後ろの$\sin x$を積分すると，微分された$\sin^{n-1}x$から$\cos x$が現れ，$\cos^2x$を使う形になる。

$n\geqq2$として部分積分を行うと
\begin{align*}
 I_n
 &=\left[-\sin^{n-1}x\cos x\right]_0^{\frac{\pi}{2}}
 +(n-1)\int_0^{\frac{\pi}{2}}\sin^{n-2}x\cos^2x\dd x\\
 &=(n-1)\int_0^{\frac{\pi}{2}}\sin^{n-2}x(1-\sin^2x)\dd x\\
 &=(n-1)(I_{n-2}-I_n).
\end{align*}
両端の項は$0$である。最後の式の$I_n$を左辺に移すと
\begin{equation}\label{式:ウォリス漸化式}
 I_n=\frac{n-1}{n}I_{n-2}\qquad(n\geqq2)
\end{equation}
を得る。指数が$2$ずつ下がるので，偶数の場合は$I_0$へ，奇数の場合は$I_1$へつながる。

\begin{theorem}{ウォリスの積分公式}
$m\geqq1$に対して
\begin{align}
 I_{2m}
 &=\frac{1\cdot3\cdot5\cdots(2m-1)}{2\cdot4\cdot6\cdots2m}
   \frac{\pi}{2},\label{式:ウォリス偶数}\\
 I_{2m+1}
 &=\frac{2\cdot4\cdot6\cdots2m}{3\cdot5\cdot7\cdots(2m+1)}.
 \label{式:ウォリス奇数}
\end{align}
\end{theorem}

例えば$I_6=\frac56I_4=\frac56\frac34I_2=\frac56\frac34\frac12I_0$であり，一般の場合もまったく同じように式\eqref{式:ウォリス漸化式}を繰り返せばよい。偶数の積と奇数の積を短く書くために
\[
 (2m)!!=2\cdot4\cdots2m,\qquad
 (2m-1)!!=1\cdot3\cdots(2m-1)
\]
と定義することがある。これは\textbf{二重階乗}と呼ばれ，階乗を二回行うという意味ではない。この記号なら
\[
 I_{2m}=\frac{(2m-1)!!}{(2m)!!}\frac{\pi}{2},\qquad
 I_{2m+1}=\frac{(2m)!!}{(2m+1)!!}
\]
となる。

\begin{example}{偶数と奇数を分けて計算する}
次の二つの積分を求めよ。
\[
 \int_0^{\frac{\pi}{2}}\sin^6x\dd x,\qquad
 \int_0^{\frac{\pi}{2}}\sin^7x\dd x.
\]
\end{example}

偶数の場合は$I_0=\frac{\pi}{2}$，奇数の場合は$I_1=1$へ戻るので
\begin{align*}
 I_6&=\frac56\frac34\frac12\frac{\pi}{2}
     =\frac{5\pi}{32},\\
 I_7&=\frac67\frac45\frac23\cdot1
     =\frac{16}{35}.
\end{align*}
偶数と奇数で最後の値が違うことが，答えに$\pi$が現れるかどうかの違いになっている。積分区間が$\left[0,\frac{\pi}{2}\right]$であることも大切である。任意の区間では，部分積分の端点の項が同じように消えるとは限らない。

\subsection{隣り合う積分の比をはさむ}

公式は積分の値を積で表したものである。今度は，積分の大小を使ってその積の極限を求める。$0\leqq\sin x\leqq1$だから，指数を大きくすると値は小さくなる。したがって
\[
 I_{n+2}\leqq I_{n+1}\leqq I_n.
\]
$I_n>0$で割り，式\eqref{式:ウォリス漸化式}を使うと
\[
 \frac{n+1}{n+2}
 =\frac{I_{n+2}}{I_n}
 \leqq\frac{I_{n+1}}{I_n}\leqq1.
\]
両端が$1$へ近づくため
\begin{equation}\label{式:ウォリス隣接比}
 \lim_{n\to\infty}\frac{I_{n+1}}{I_n}=1
\end{equation}
である。各$I_n$の極限を先に求めなくても，比の極限は分かる。

ここで，次の有限の積を作る。
\[
 P_m=\prod_{k=1}^{m}\frac{(2k)^2}{(2k-1)(2k+1)}.
\]
$\prod$は，指定された範囲の項をすべて掛ける記号である。最初の数項を書けば
\[
 P_m=\frac{2\cdot2}{1\cdot3}
      \frac{4\cdot4}{3\cdot5}\cdots
      \frac{2m\cdot2m}{(2m-1)(2m+1)}.
\]
偶数と奇数の積をそれぞれまとめ，式\eqref{式:ウォリス偶数}と式\eqref{式:ウォリス奇数}を使うと
\begin{align*}
 \frac{I_{2m+1}}{I_{2m}}
 &=\frac{2}{\pi}
   \frac{\{(2m)!!\}^2}{(2m-1)!!(2m+1)!!}\\
 &=\frac{2}{\pi}P_m.
\end{align*}
左辺は$1$へ近づくので，次の極限が得られる。
\begin{equation}\label{式:ウォリスの積}
 \boxed{\displaystyle
 \lim_{m\to\infty}\prod_{k=1}^{m}
 \frac{4k^2}{4k^2-1}=\frac{\pi}{2}}
\end{equation}
これが\textbf{ウォリスの積}である。無限個の数を一度に掛けたと考えるのではなく，有限の積$P_m$の数列の極限として定義している。

証明の途中の不等式を残しておけば，近似の誤差も評価できる。$n=2m$の場合には
\[
 \frac{2m+1}{2m+2}
 \leqq\frac{I_{2m+1}}{I_{2m}}\leqq1
\]
だから
\[
 \frac{2m+1}{2m+2}\frac{\pi}{2}
 \leqq P_m\leqq\frac{\pi}{2}.
\]
従って，正の数を掛けて整理すれば
\[
 2P_m\leqq\pi\leqq\frac{2m+2}{2m+1}\,2P_m.
\]
公式が極限を与えるだけでなく，有限の計算の結果がどれほど確かな範囲を与えるかも分かる。

%
\subsection{中央の二項係数との関係}

偶数公式の分数には，階乗を使った別の形がある。
\[
 (2m)!!=2^m m!,\qquad
 (2m-1)!!=\frac{(2m)!}{2^m m!}
\]
だから
\begin{equation}\label{式:ウォリスと二項係数}
 I_{2m}=\frac{\pi}{2}\frac{(2m)!}{2^{2m}(m!)^2}
 =\frac{\pi}{2}\frac{\binom{2m}{m}}{4^m}.
\end{equation}
ここで$\binom{2m}{m}=\frac{(2m)!}{(m!)^2}$は二項係数である。

さらに偶数公式と奇数公式を掛けると，多くの因子が消えて
\[
 I_{2m}I_{2m+1}=\frac{\pi}{2(2m+1)}
\]
となる。これと$\frac{I_{2m+1}}{I_{2m}}\to1$を組み合わせれば
\[
 mI_{2m}^2
 =\frac{m\pi}{2(2m+1)}\frac{I_{2m}}{I_{2m+1}}
 \longrightarrow\frac{\pi}{4}.
\]
$I_{2m}>0$なので平方根を取ることができ，
\[
 \sqrt m\,I_{2m}\longrightarrow\frac{\sqrt\pi}{2}.
\]
式\eqref{式:ウォリスと二項係数}に戻すと
\begin{equation}\label{式:中央二項係数の極限}
 \lim_{m\to\infty}\frac{\sqrt m}{4^m}\binom{2m}{m}
 =\frac1{\sqrt\pi}
\end{equation}
を得る。階乗を含む数列の大きさが，三角関数の積分を経由して分かったことになる。

%
%
\section{多項式の積分からベータ積分へ}\label{節:ベータ積分}

\subsection{非負整数の指数}

非負整数$m,n$に対して
\[
 J_{m,n}=\int_0^1x^m(1-x)^n\dd x
\]
を考える。$m,n$が小さければ展開で求められる。しかし指数が大きくなると，展開した項数も多くなる。ここでは$x^m$と$(1-x)^n$をそのまま残し，部分積分で指数を一つずつ変える。

$n\geqq1$として，$x^m$を積分し，$(1-x)^n$を微分すると
\begin{align*}
 J_{m,n}
 &=\left[\frac{x^{m+1}}{m+1}(1-x)^n\right]_0^1
   +\frac{n}{m+1}\int_0^1x^{m+1}(1-x)^{n-1}\dd x\\
 &=\frac{n}{m+1}J_{m+1,n-1}.
\end{align*}
両端で積が$0$になるので，端点の項は消える。この式を$n$回繰り返すと
\[
 J_{m,n}
 =\frac{n!}{(m+1)(m+2)\cdots(m+n)}J_{m+n,0}.
\]
最後の積分は
\[
 J_{m+n,0}=\int_0^1x^{m+n}\dd x=\frac1{m+n+1}
\]
だから
\begin{equation}\label{式:整数ベータ積分}
 \boxed{\displaystyle
 \int_0^1x^m(1-x)^n\dd x=\frac{m!n!}{(m+n+1)!}}
\end{equation}
となる。$n=0$の場合も，右辺を直接整理すれば同じ式で表される。

\begin{example}{展開せずに多項式の積分を求める}
次の積分を求めよ。
\[
 \int_0^1x^3(1-x)^2\dd x.
\]
\end{example}

式\eqref{式:整数ベータ積分}で$m=3,n=2$とすると
\[
 \int_0^1x^3(1-x)^2\dd x
 =\frac{3!2!}{6!}=\frac1{60}.
\]
展開で検算すると
\[
 \int_0^1(x^3-2x^4+x^5)\dd x
 =\frac14-\frac25+\frac16=\frac1{60}
\]
となり，一致する。公式は展開を不要にするが，指数がどこに入るかを確かめるため，小さい場合を計算してみることにも意味がある。

%
\subsection{端点で大きくなる関数の積分}

指数を正の実数へ広げるため，$p,q>0$として
\[
 x^{p-1}(1-x)^{q-1}
\]
を考えたい。例えば$p=\frac12$なら$x^{p-1}=\frac1{\sqrt x}$となり，$x=0$で値が定まらない。閉区間で連続な関数の積分という，これまでの定義をそのまま使うことはできない。

そこで，問題のある端点をいったん避けて積分し，その端点へ近づける。$0<c<1$を一つ固定して
\begin{equation}\label{式:両端広義積分の定義}
 \int_0^1 f(x)\dd x
 =\lim_{\varepsilon\downarrow0}\int_{\varepsilon}^{c}f(x)\dd x
  +\lim_{\delta\downarrow0}\int_c^{1-\delta}f(x)\dd x
\end{equation}
と定める。ただし，右辺の\emph{二つの極限がそれぞれ有限の値を持つ場合}に限る。このような積分を\textbf{広義積分}といい，有限の値を持つことを収束するという。$c$を別の内点に変えても，その二点間の通常の積分が一方から他方へ移るだけなので，和の値は変わらない。

両端へ同じ速さで近づける$\int_{\varepsilon}^{1-\varepsilon}$の極限だけでは，この定義を満たしたことにはならない。一方の端で正に発散し，他方の端で負に発散するものが打ち消し合うおそれがあるためである。この章のベータ積分は非負関数だが，後にも使えるよう，両端を独立に調べる定義にしておく。

まず，$p>0$なら
\[
 \int_{\varepsilon}^{c}x^{p-1}\dd x
 =\frac{c^p-\varepsilon^p}{p}
 \longrightarrow\frac{c^p}{p}
\]
である。$p$が$1$より小さくても，$\varepsilon^p\to0$となるので収束する。関数の高さが無限に大きくなることと，面積が無限に大きくなることは同じではない。

ベータ積分の左半分では，$x\in\left[0,\frac12\right]$において$(1-x)^{q-1}$が連続で，ある正の定数$M$以下である。従って
\[
 0\leqq x^{p-1}(1-x)^{q-1}\leqq Mx^{p-1}.
\]
下端$\varepsilon$を小さくすると積分値は増加するが，上の評価により$\frac{M}{p}\left(\frac12\right)^p$以下にとどまる。単調で上に有界な数列の極限と同じ考え方で，この極限は有限に存在する。右半分も$u=1-x$と置けば$q>0$によって同様に収束する。したがって$p,q>0$なら，両端の広義積分がともに定まる。

\begin{definition}{ベータ関数}
$p,q>0$に対して
\begin{equation}\label{式:ベータ関数の定義}
 B(p,q)=\int_0^1x^{p-1}(1-x)^{q-1}\dd x
\end{equation}
と定める。端点で関数が定まらない場合には，式\eqref{式:両端広義積分の定義}による広義積分とする。
\end{definition}

整数の場合との対応は，指数に$1$のずれがある。式\eqref{式:整数ベータ積分}は
\[
 B(m+1,n+1)=\frac{m!n!}{(m+n+1)!}
\]
と書ける。$B(m,n)$の指数は$m-1,n-1$であり，$m,n$そのものではない。記号の引数と被積分関数の指数を取り違えないようにする。

%
\subsection{対称性と漸化式}

$u=1-x$と置くと二つの指数が入れ替わるため
\begin{equation}\label{式:ベータ対称性}
 B(p,q)=B(q,p)
\end{equation}
である。広義積分の場合も，端点を避けた区間で置換し，各端点の極限を取れば同じ結論になる。

また$x+(1-x)=1$を掛ければ
\begin{equation}\label{式:ベータ和の関係}
 B(p+1,q)+B(p,q+1)=B(p,q).
\end{equation}
もう一つの関係を得るために$x^p(1-x)^q$を微分する。
\[
 \frac{\dd}{\dd x}\{x^p(1-x)^q\}
 =p x^{p-1}(1-x)^q-qx^p(1-x)^{q-1}.
\]
これを$\varepsilon$から$1-\delta$まで積分し，両端をそれぞれ極限へ近づける。$p,q>0$なので$x^p(1-x)^q$は両端で$0$へ近づき，右辺の二つの積分はすでに収束すると分かっている。よって
\[
 pB(p,q+1)=qB(p+1,q).
\]
式\eqref{式:ベータ和の関係}と連立すると
\begin{equation}\label{式:ベータ漸化式}
 B(p+1,q)=\frac{p}{p+q}B(p,q),\qquad
 B(p,q+1)=\frac{q}{p+q}B(p,q)
\end{equation}
を得る。この漸化式は整数以外にも成り立つ。整数のときには，$B(1,1)=1$から出発して階乗の公式を再び導くこともできる。

%
\subsection{三角関数の積分とのつながり}

$x=\sin^2t$と置くと，$0<t<\frac{\pi}{2}$で$\dd x=2\sin t\cos t\dd t$である。$\sin t,\cos t$が正なので，実数の指数もそのまま扱える。式\eqref{式:ベータ関数の定義}は
\begin{equation}\label{式:ベータ三角表示}
 B(p,q)=2\int_0^{\frac{\pi}{2}}
 (\sin t)^{2p-1}(\cos t)^{2q-1}\dd t
\end{equation}
となる。厳密には端点を避けて置換した後で極限を取る。左辺の両端の積分が収束しているので，対応する右辺の極限も有限である。

$p=q=\frac12$なら両方の指数が$0$になり，
\[
 B\left(\frac12,\frac12\right)
 =2\int_0^{\frac{\pi}{2}}1\dd t=\pi.
\]
つまり
\[
 \int_0^1\frac1{\sqrt{x(1-x)}}\dd x=\pi
\]
である。両端で高さが限りなく大きくなるにもかかわらず，積分は有限になる具体例である。

一方，$p=\frac{n+1}{2}$，$q=\frac12$とすれば
\[
 B\left(\frac{n+1}{2},\frac12\right)=2I_n.
\]
ベータ関数の漸化式で$p$を$1$増やすことは，$\sin$の指数$n$を$2$増やすことに対応する。実際
\[
 B\left(\frac{n+3}{2},\frac12\right)
 =\frac{n+1}{n+2}
  B\left(\frac{n+1}{2},\frac12\right)
\]
は$I_{n+2}=\frac{n+1}{n+2}I_n$そのものである。整数の多項式積分とウォリスの積分は，ベータ積分の異なる場合として結び付いた。

%
%
\section*{章末問題}
\begin{enumerate}[label=\textbf{\thechapter.\arabic*},leftmargin=*]
 \item\label{問:区分求積分子が添字}
 次の極限を求めよ。
 \[
  \lim_{n\to\infty}\sum_{k=1}^{n}\frac{k}{n^2+k^2}.
 \]
 \item\label{問:区分求積根号}
 次の極限を求めよ。
 \[
  \lim_{n\to\infty}\cfrac1{n^{\cfrac32}}
      \sum_{k=1}^{n}\sqrt{n+k}.
 \]
 \item\label{問:逆平方和の残り}
 自然数$N>n\geqq1$に対して
 \[
  \frac1{n+1}-\frac1{N+1}
  \leqq\sum_{k=n+1}^{N}\frac1{k^2}
  \leqq\frac1n-\frac1N
 \]
 を示せ。これを使って$\sum_{k=1}^{\infty}\frac1{k^2}$が収束し，その第$n$項より後の和$R_n$が
 $\frac1{n+1}\leqq R_n\leqq\frac1n$を満たすことを示せ。
 \item\label{問:積分制約と二乗}
 $f$は$[0,1]$で連続で，$\int_0^1xf(x)\dd x=1$を満たすとする。
 $\int_0^1f(x)^2\dd x$の最小値を求め，それを与える関数を求めよ。
 \item\label{問:対数の積分漸化式}
 非負整数$n$に対して$K_n=\int_1^{\e}(\log x)^n\dd x$とおく。
 $K_0$と，$n\geqq1$における漸化式を求め，$K_3$を計算せよ。
 \item\label{問:ウォリス積の誤差}
 $P_m=\prod_{k=1}^{m}\frac{4k^2}{4k^2-1}$とする。
 \[
 0\leqq\frac{\pi}{2}-P_m\leqq\frac{\pi}{4m+4}
 \]
 を示せ。また$\lim_{m\to\infty}\sqrt m\int_0^{\frac{\pi}{2}}\cos^{2m}x\dd x$を求めよ。
 \item\label{問:整数ベータ二項係数}
 非負整数$m,n$に対して
 \[
 \int_0^1\binom{m+n}{m}x^m(1-x)^n\dd x=\frac1{m+n+1}
 \]
 を示せ。次に$\int_0^1x^2(1-x)^4\dd x$を求めよ。
 \item\label{問:半整数ベータ}
 広義積分$\int_0^1\sqrt{\frac{x}{1-x}}\dd x$の収束を両端で確かめ，ベータ関数の漸化式を使って値を求めよ。さらに$x=\sin^2t$による置換でも検算せよ。
\end{enumerate}

\smallskip\noindent 解答は\pageref{解答:第9章}ページ。
%
